% =====================================================================
%  Relatório Técnico
% =====================================================================
\documentclass[11pt,a4paper]{article}

% ---------- Codificação e idioma ----------
\usepackage[utf8]{inputenc}
\usepackage[T1]{fontenc}
\usepackage[brazil]{babel}

% ---------- Layout ----------
\usepackage[a4paper,margin=2.5cm]{geometry}
\usepackage{setspace}
\onehalfspacing

% ---------- Recursos ----------
\usepackage{graphicx}      % figuras
\usepackage{booktabs}      % tabelas
\usepackage{amsmath}       % fórmulas
\usepackage{xcolor}        % cores
\usepackage{listings}      % código-fonte
\usepackage[hidelinks]{hyperref}

% ---------- Estilo para código C ----------
\lstdefinestyle{c}{
    language=C,
    basicstyle=\ttfamily\small,
    keywordstyle=\color{blue},
    commentstyle=\color{gray},
    stringstyle=\color{teal},
    numbers=left,
    numberstyle=\tiny\color{gray},
    breaklines=true,
    frame=single,
    tabsize=4,
    showstringspaces=false
}
\lstset{style=c}

% ---------- Metadados ----------
\title{Projetando um Gerenciador de Memória Virtual\\
       \large Trabalho Prático 2 - Sistemas Operacionais}
\author{
    Davi Lage \\
    Gabriel Santos \\
    Isaías Alves \\
    Laura Noronha \\
    Lucas Spiazzi
}
\date{\today}

% =====================================================================
\begin{document}

\maketitle

\begin{abstract}
% TODO: Resumir, em 5–8 linhas, o objetivo do trabalho, o que foi
% implementado e os principais resultados obtidos.
\end{abstract}

\tableofcontents
\newpage

% =====================================================================
\section{Introdução e Fundamentação Teórica}
\label{sec:introducao}
% TODO: Apresentar o tema (memória virtual e tradução de endereços) e a
% fundamentação teórica (paginação, tabela de páginas, TLB, page fault,
% paginação por demanda e substituição de páginas por LRU/Aging). Encerrar
% com o objetivo do trabalho e a organização do relatório.

% =====================================================================
\section{Arquitetura e Estruturas de Dados}
\label{sec:arquitetura}
% TODO: Descrever a arquitetura modular do simulador (papel de cada módulo
% e a separação include/ vs src/) e as estruturas de dados utilizadas,
% explicando a função de cada uma. Ilustrar o fluxo de tradução.

% =====================================================================
\section{Implementação e Decisões de Projeto}
\label{sec:implementacao}
% TODO: Descrever como cada parte foi implementada e justificar as decisões
% de projeto: extração de página/offset, paginação por demanda, gerência de
% quadros, substituição por LRU aproximado (Aging), coerência na expulsão,
% TLB com FIFO, cálculo do endereço físico e estatísticas.

% =====================================================================
\section{Testes Realizados e Análise dos Resultados}
\label{sec:resultados}
% TODO: Descrever a metodologia de testes e apresentar/analisar os resultados
% obtidos com os arquivos de entrada (random e location), comparando as taxas
% de page fault e de acerto do TLB.

% =====================================================================
\section{Conclusão}
\label{sec:conclusao}
% TODO: Sintetizar os objetivos alcançados, as principais decisões e
% aprendizados, e a relação entre a teoria e a implementação.

% =====================================================================
\section{Uso de Inteligência Artificial}
\label{sec:ia}
% TODO: Documentar como a IA foi utilizada (geração de código, explicações,
% depuração, testes), listar TODOS os prompts empregados e discutir a
% metodologia Spec-Driven Development (SDD).

% =====================================================================
\section*{Repositório}
% TODO: Inserir o link do repositório público no GitHub.

% =====================================================================
% TODO: Incluir as referências bibliográficas utilizadas.

\end{document}
